\documentclass[10pt,a4paper]{amsart}
\usepackage[margin=19mm]{geometry}
\usepackage{amsmath,amssymb,mathtools}
\usepackage[scr=rsfs,cal=euler]{mathalfa}
\usepackage{xfrac}
\usepackage[unicode,hidelinks,bookmarks=false]{hyperref}
\newcommand{\ie}{i.e.\ }
\newcommand{\cf}{cf.\ }
\newcommand{\resp}{respectively\ }
\newcommand{\Subsection}{\subsection}
\providecommand{\nobreakdash}{\nobreak\hbox{-}}
\setlength{\emergencystretch}{2em}
\multlinegap0pt
\usepackage[normalem]{ulem}
\usepackage{amssymb}
\usepackage{mathtools}
\usepackage{xcolor}
\usepackage{bm}
\newtheorem{lemma}{Lemma}
\title{Self-Similar Singular Solutions to the Nonlinear Schrödinger and the Complex Ginzburg-Landau Equations}
\author{Joel Dahne, Jordi-Llu\'is Figueras}
\date{Source: arXiv:2410.05480v3 (CC BY 4.0)}
\begin{document}
\maketitle

\noindent\emph{Source and licence.} Self-Similar Singular Solutions to the Nonlinear Schrödinger and the Complex Ginzburg-Landau Equations. Version arXiv:2410.05480v3 (CC BY 4.0), distributed by arXiv under the Creative Commons Attribution 4.0 International licence, http://creativecommons.org/licenses/by/4.0/, as recorded in the arXiv licence metadata for this exact version and shown on the abstract page https://arxiv.org/abs/2410.05480v3. Full licence notice: \texttt{licenses/arxiv-2410.05480-CC-BY-4.0.txt}.\par\smallskip

\noindent\textit{Editorial scope.} Counted unit: the article's lemma lemma:tail-bound (source0.tex lines 4314--4341). The article's complete original proof of it is delivered with it (source0.tex lines 4343--4456, one \texttt{proof} environment of 114 lines). No mathematics is written, repaired or re-derived: the statement and the proof are the article's own lines, in the article's own notation, and no retained line is substituted or re-flowed.  No further source-local context is needed: every cross-reference of every retained line resolves inside the two delivered spans.  Nothing was omitted from any retained span, so the package carries no omission record.  No retained line contains a citation, so the wrapper carries no bibliography and no bibliography entry.  No retained line contains a figure environment, an image-inclusion command or drawing code, so no image is delivered and the package declares no asset dependency.  Editorial formatting only, in addition: an AMS \texttt{amsart} preamble that restates the article's own declarations and macros used by the retained lines, namely the environments \texttt{lemma} on separate counters, together with the standard packages those lines need.  The retained environments print in this document as Lemma 1 in order of appearance, whereas the article prints the same items with the numbering of its own full document.  The complete native source of the article as distributed in the arXiv e-print for this exact version is bundled unchanged as \texttt{sources/166459/source0.tex}.
\begin{lemma}\label{lemma:tail-bound}
  Let \(\sigma = 1\). Let \(M\), \(N\), \(C\) and \(r\) be such that
  \(N\) is even, \(3M < N\),
  \begin{equation*}
    |a_n|, |b_n| \leq C r^n \quad\text{for}\quad n < M
  \end{equation*}
  and
  \begin{equation*}
    |a_n|, |b_n| \leq r^n \quad\text{for}\quad M \leq n \leq N.
  \end{equation*}
  If
  \begin{multline}
    \label{eq:tail-bound-inequality}
    \frac{(1 + |\epsilon|)}{(1 + \epsilon^{2})}\Bigg(
      \frac{|\kappa|}{N + d}
      + \frac{|\omega|}{(N + 2)(N + d)}\\
      + 2(1 + |\delta|)\left(
        \frac{1}{8} + \frac{1}{2N}
        + \frac{3mC(1 + 3/N)}{2(N + d)}
        + \frac{3m^{2}C^{2}}{(N + 2)(N + d)}
      \right)
    \Bigg) \leq r^2,
  \end{multline}
  where \(m = \lceil M / 2 \rceil\), then
  \begin{equation*}
    |a_{n}|, |b_{n}| \leq r^{n} \quad\text{for}\quad n > N.
  \end{equation*}
\end{lemma}

\begin{proof}
  It is enough to show that
  \begin{equation*}
    |a_{N + 1}|, |b_{N + 1}| \leq r^{N + 1},
    |a_{N + 2}|, |b_{N + 2}| \leq r^{N + 2}
  \end{equation*}
  and use that the left-hand side in \eqref{eq:tail-bound-inequality}
  is decreasing in \(N\). The result then follows for \(n > N + 2\) by
  induction.

  Since \(a_{n}\) and \(b_{n}\) are zero for odd \(n\) and \(N\) is
  even it follows that \(a_{N + 1}\) and \(b_{N + 1}\) are zero. We
  therefore only have to prove that
  \(|a_{N + 2}|, |b_{N + 2}| \leq r^{N + 2}\). We give the proof for
  \(a_{N + 2}\), the bound for \(b_{N + 2}\) follows in exactly the
  same way.

  With \(\sigma = 1\) we get \(u_{1,N} = (a^{3})_N + (ab^{2})_N\) and
  \(u_{2,N} = (a^{2}b)_N + (b^{3})_N\). We have
  \begin{equation*}
    (a^{3})_{N} = \sum_{i + j + k = N}a_{i}a_{j}a_{k}.
  \end{equation*}
  However, since \(a_{n} = 0\) for odd \(n\) and \(N\) is even we can
  limit the sum to
  \begin{equation*}
    (a^{3})_{N} = \sum_{i + j + k = N/2}a_{2i}a_{2j}a_{2k}.
  \end{equation*}
  Since the bound for \(a_{n}\) depends on if \(n\) is less than \(M\)
  or not, let us group the terms in the sum by how many of the indices
  \(i, j, k\) are less than \(m = \lceil M / 2 \rceil\):
  \begin{enumerate}
  \item The number of terms with no indices less than \(m\) is
    \begin{equation*}
      T_{0} = \sum_{i = 0}^{N/2 - 3m}\sum_{j = 0}^{N/2 - 3m - i} 1 = \frac{(N/2 - 3m + 1)(N/2 - 3m + 2)}{2}.
    \end{equation*}
    These are bounded by \(r^{N}\).
  \item For the terms with exactly one index less than \(m\) there are
    \(3\) choices for which index is the small one. If we let \(i\) be
    the small index, then there are \(N / 2 - 2m - i + 1\) choices for
    \(j\) and \(k\). Summing over all choices for \(i\) and
    multiplying by \(3\) we get
    \begin{equation*}
      T_{1} = 3\sum_{i = 0}^{m - 1} (\frac{N}{2} - 2m - i + 1) = 3\frac{m}{2}\left(N - 5m + 3\right).
    \end{equation*}
    These are bounded by \(Cr^{N}\).
  \item The number of terms with exactly two indices less than \(m\)
    is \(T_{2} = 3m^{2}\). These are bounded by \(C^{2}r^{N}\).
  \end{enumerate}
  Note that since \(3M < N\) there are no terms with three indices
  less than \(m\). In total this gives us the bound
  \begin{equation*}
    |(a^{3})_{N}| \leq (T_{0} + T_{1}C + T_{2}C^{2})r^{N}.
  \end{equation*}

  Since the bounds for \(|a_n|\) and \(|b_n|\) are the same we get the
  same bound for \(|(ab^2)_N|\), \(|(a^2b)_N|\) and \(|(b^3)_N|\). In
  particular, this gives us
  \begin{equation*}
    |u_{1,N}|, |u_{2,N}| \leq 2(T_{0} + T_{1}C + T_{2}C^{2})r^{N}.
  \end{equation*}

  Combining the above bounds for \(|u_{1,N}|\) and \(|u_{2,N}|\) with
  the bounds for \(|a_n|\) and \(|b_n|\) we get
  \begin{equation*}
    |F_{1,N}| \leq (|\kappa|(N + 1) + |\omega|)r^{N} + 2(1 + |\delta|)(T_{0} + T_{1}C + T_{2}C^{2})r^{N},
  \end{equation*}
  and the same for \(|F_{2,N}|\). This gives us
  \begin{equation*}
    \begin{split}
      |a_{N + 2}|
      &\leq \frac{(1 + |\epsilon|)(|\kappa|(N + 1) + |\omega| + 2(1 + |\delta|)(T_{0} + T_{1}C + T_{2}C^{2}))}{(N + 2)(N + d)(1 + \epsilon^{2})}r^{N}\\
      &= \frac{(1 + |\epsilon|)}{(1 + \epsilon^{2})}\left(
        \frac{|\kappa|}{N + d}
        + \frac{|\omega|}{(N + 2)(N + d)}
        + 2(1 + |\delta|)\frac{(T_{0} + T_{1}C + T_{2}C^{2})}{(N + 2)(N + d)}
        \right)r^{N}.
    \end{split}
  \end{equation*}
  We have
  \begin{equation*}
    \begin{split}
      \frac{T_{0}}{(N + 2)(N + d)}
      &= \frac{(N/2 - 3m + 1)(N/2 - 3m + 2)}{2(N + 2)(N + d)}\\
      &= \frac{1}{8}\frac{(N - 6m + 2)(N - 6m + 4)}{(N + 2)(N + d)}\\
      &\leq \frac{1}{8}\frac{N + 4}{N + d}
        \leq \frac{1}{8}\left(1 + \frac{4}{N}\right)
        = \frac{1}{8} + \frac{1}{2N},
    \end{split}
  \end{equation*}
  as well as
  \begin{equation*}
    \frac{T_{1}}{(N + 2)(N + d)}
    = 3\frac{m}{2}\frac{N - 5m + 3}{(N + 2)(N + d)}
    \leq 3\frac{m}{2}\frac{N + 3}{(N + 2)(N + d)}
    \leq \frac{3m(1 + 3/N)}{2(N + d)}.
  \end{equation*}
  This gives us
  \begin{multline*}
    |a_{N + 2}|
    \leq \frac{(1 + |\epsilon|)}{(1 + \epsilon^{2})}\Bigg(
    \frac{|\kappa|}{N + d}
      + \frac{|\omega|}{(N + 2)(N + d)}\\
      + 2(1 + |\delta|)\left(
        \frac{1}{8} + \frac{1}{2N}
        + \frac{3mC(1 + 3/N)}{2(N + d)}
        + \frac{3m^{2}C^{2}}{(N + 2)(N + d)}
      \right)
    \Bigg)r^{N}.
  \end{multline*}

  By~\eqref{eq:tail-bound-inequality}, the factor in front of
  \(r^{N}\) is bounded by \(r^{2}\), giving us our required bound
  \(|a_{N + 2}| \leq r^{N + 2}\).
\end{proof}

\end{document}
